\documentclass[12pt,letterpaper]{article}

% ── Encoding & fonts ──────────────────────────────────────────────────────────
\usepackage[T1]{fontenc}
\usepackage[utf8]{inputenc}
\usepackage{lmodern}

% ── Page geometry ─────────────────────────────────────────────────────────────
\usepackage[top=2.5cm, bottom=2.5cm, left=2.5cm, right=2.5cm]{geometry}
\setlength{\headheight}{15pt}

% ── Math ──────────────────────────────────────────────────────────────────────
\usepackage{amsmath, amssymb, amsthm}
\usepackage{mathtools}
\usepackage{bm}

% ── Lists ─────────────────────────────────────────────────────────────────────
\usepackage{enumitem}

% ── Tables ────────────────────────────────────────────────────────────────────
\usepackage{booktabs}
\usepackage{array}

% ── Graphics / colour ─────────────────────────────────────────────────────────
\usepackage{xcolor}
\usepackage{graphicx}

% ── Coloured answer boxes ─────────────────────────────────────────────────────
\usepackage[most]{tcolorbox}

% ── Headers / footers ─────────────────────────────────────────────────────────
\usepackage{fancyhdr}
\usepackage{lastpage}

% ── Section formatting ────────────────────────────────────────────────────────
\usepackage{titlesec}
\usepackage{float}
\usepackage{microtype}

% ── Hyperlinks & PDF metadata ─────────────────────────────────────────────────
\usepackage[colorlinks=true, linkcolor=TorontoBlue, urlcolor=TorontoBlue,
            citecolor=TorontoBlue, pdftitle={CPCS110 A6 Solution Key},
            pdfauthor={TorontoMU Physics}]{hyperref}

% ── parskip loaded AFTER hyperref to avoid \@setpar conflicts ─────────────────
\usepackage{parskip}

% ═══════════════════════════════════════════════════════════════════════════════
%  Colour palette
% ═══════════════════════════════════════════════════════════════════════════════
\definecolor{TorontoBlue}{RGB}{0, 52, 120}
\definecolor{TorontoGold}{RGB}{253, 182, 0}
\definecolor{AnswerBg}{RGB}{235, 244, 255}
\definecolor{AnswerBorder}{RGB}{0, 52, 120}
\definecolor{GivenBg}{RGB}{245, 250, 235}
\definecolor{GivenBorder}{RGB}{70, 130, 40}
\definecolor{NoteBg}{RGB}{255, 250, 230}
\definecolor{NoteBorder}{RGB}{200, 150, 0}
\definecolor{CheckBg}{RGB}{240, 255, 240}
\definecolor{CheckBorder}{RGB}{0, 140, 0}

% Predefine tinted named colours for cross-distribution robustness
\colorlet{TorontoBlue80}{TorontoBlue!80}
\colorlet{TorontoBlue60}{TorontoBlue!60}
\colorlet{GivenTitleBg}{GivenBorder!50}

% ═══════════════════════════════════════════════════════════════════════════════
%  tcolorbox styles
% ═══════════════════════════════════════════════════════════════════════════════
\tcbset{
  answerbox/.style={
    enhanced, breakable,
    colback=AnswerBg, colframe=AnswerBorder,
    boxrule=1.5pt, arc=4pt,
    fonttitle=\bfseries\color{white},
    coltitle=white,
    attach boxed title to top left={yshift=-2mm, xshift=6mm},
    boxed title style={colback=AnswerBorder, arc=3pt, boxrule=0pt},
    title={Final Answer},
    top=4pt, bottom=4pt, left=6pt, right=6pt
  },
  givenbox/.style={
    enhanced, breakable,
    colback=GivenBg, colframe=GivenBorder,
    boxrule=1pt, arc=4pt,
    fonttitle=\bfseries\small,
    coltitle=black,
    attach boxed title to top left={yshift=-2mm, xshift=6mm},
    boxed title style={colback=GivenTitleBg, arc=3pt, boxrule=0pt,
                       coltext=black},
    title={Given Quantities},
    top=4pt, bottom=4pt, left=6pt, right=6pt
  },
  checkbox/.style={
    enhanced,
    colback=CheckBg, colframe=CheckBorder,
    boxrule=1pt, arc=4pt,
    top=3pt, bottom=3pt, left=6pt, right=6pt
  },
  principbox/.style={
    enhanced,
    colback=NoteBg, colframe=NoteBorder,
    boxrule=1pt, arc=4pt,
    fonttitle=\bfseries,
    coltitle=black,
    attach boxed title to top left={yshift=-2mm, xshift=6mm},
    boxed title style={colback=NoteBorder!40, arc=3pt, boxrule=0pt,
                       coltext=black},
    title={Physics Principle},
    top=4pt, bottom=4pt, left=6pt, right=6pt
  }
}

% ═══════════════════════════════════════════════════════════════════════════════
%  Header / footer
% ═══════════════════════════════════════════════════════════════════════════════
\pagestyle{fancy}
\fancyhf{}
\fancyhead[L]{\small\color{TorontoBlue}\textbf{CPCS110 --- Matter \& Interactions}}
\fancyhead[R]{\small\color{TorontoBlue}\textbf{Assignment 6 Solution Key}}
\fancyfoot[C]{\small\color{gray}Page \thepage\ of \pageref{LastPage}}
\renewcommand{\headrulewidth}{1.5pt}
\renewcommand{\footrulewidth}{0.4pt}
\renewcommand{\headrule}{\color{TorontoBlue}\hrule width\headwidth height\headrulewidth}

\fancypagestyle{plain}{
  \fancyhf{}
  \fancyfoot[C]{\small\color{gray}Page \thepage\ of \pageref{LastPage}}
  \renewcommand{\headrulewidth}{0pt}
  \renewcommand{\footrulewidth}{0.4pt}
}

% ═══════════════════════════════════════════════════════════════════════════════
%  Section formatting  (using predefined tinted colours)
% ═══════════════════════════════════════════════════════════════════════════════
\titleformat{\section}[block]
  {\normalfont\Large\bfseries\color{TorontoBlue}}
  {\thesection.}{0.8em}{}
  [\color{TorontoBlue}\titlerule]

\titleformat{\subsection}[block]
  {\normalfont\large\bfseries\color{TorontoBlue80}}
  {\thesubsection}{0.6em}{}

\titleformat{\subsubsection}[block]
  {\normalfont\normalsize\bfseries\color{TorontoBlue60}}
  {\thesubsubsection}{0.5em}{}

% ═══════════════════════════════════════════════════════════════════════════════
%  Convenience macros
% ═══════════════════════════════════════════════════════════════════════════════
\newcommand{\vv}[1]{\vec{#1}}
\newcommand{\bvec}[1]{\langle #1 \rangle}   % Matter & Interactions component notation
\newcommand{\units}[1]{\;\text{#1}}
\newcommand{\sci}[2]{#1 \times 10^{#2}}
\newcommand{\pts}[1]{\hfill\textcolor{TorontoBlue60}{[#1 pts]}}

% ═══════════════════════════════════════════════════════════════════════════════
\begin{document}
% ═══════════════════════════════════════════════════════════════════════════════

% ── Title page ────────────────────────────────────────────────────────────────
\begin{titlepage}
\thispagestyle{plain}
\begin{center}
  \vspace*{1.5cm}

  {\color{TorontoBlue}\rule{\linewidth}{2pt}}\\[0.6cm]

  {\Huge\bfseries\color{TorontoBlue} CPCS110 --- Matter \& Interactions}\\[0.4cm]
  {\Large\bfseries\color{TorontoGold} Assignment 6 Solution Key}\\[0.3cm]
  {\large\color{gray} Chapters 5--6: Potential Energy \& Electric Potential}\\[0.4cm]

  {\color{TorontoBlue}\rule{\linewidth}{2pt}}\\[0.8cm]

  \begin{tabular}{ll}
    \toprule
    \textbf{Course}   & CPCS110 --- Matter \& Interactions \\[4pt]
    \textbf{Term}     & Fall 2025 \\[4pt]
    \textbf{Topics}   & Chapters 5--6 \\[4pt]
    \textbf{Status}   & \textcolor{CheckBorder}{\bfseries Verified Solution Key} \\[4pt]
    \bottomrule
  \end{tabular}

  \vspace{1.2cm}

  \begin{tcolorbox}[
    enhanced, width=0.85\linewidth,
    colback=NoteBg, colframe=NoteBorder,
    boxrule=1pt, arc=6pt,
    fonttitle=\bfseries, title={Notation Guide},
    attach boxed title to top left={yshift=-2mm, xshift=6mm},
    boxed title style={colback=NoteBorder!40, arc=3pt, boxrule=0pt}
  ]
  \textbf{Vector notation:} This solution key uses the \emph{Matter \& Interactions}
  component notation $\bvec{A_x, A_y, A_z}$ (angle brackets) to represent vectors.
  This is equivalent to $A_x\,\hat{x} + A_y\,\hat{y} + A_z\,\hat{z}$.\\[4pt]
  \textbf{Significant figures:} Answers are reported to 4 significant figures unless
  the given data suggest fewer.\\[4pt]
  \textbf{Constants used:} $k = 8.988\times10^{9}\units{N\cdot m^{2}/C^{2}}$,
  $e = 1.6\times10^{-19}\units{C}$, $m_p = 1.673\times10^{-27}\units{kg}$,
  $m_e = 9.109\times10^{-31}\units{kg}$.
  \end{tcolorbox}

  \vspace{0.8cm}

  {\bfseries\color{TorontoBlue} Topics Covered}\\[0.4cm]
  \begin{minipage}{0.85\linewidth}
  \begin{itemize}[leftmargin=1.5em, itemsep=4pt]
    \item[\textcolor{TorontoBlue}{$\triangleright$}] \textbf{[Q1--Q2]} Gravitational potential energy and orbital energy
    \item[\textcolor{TorontoBlue}{$\triangleright$}] \textbf{[Q3--Q4]} Gravitational PE and escape speed
    \item[\textcolor{TorontoBlue}{$\triangleright$}] \textbf{[Q5]} Electric potential energy of point charges
    \item[\textcolor{TorontoBlue}{$\triangleright$}] \textbf{[Q6]} Electric field from potential: $\vec{E} = -\nabla V$
    \item[\textcolor{TorontoBlue}{$\triangleright$}] \textbf{[Q7]} Coulomb field: unit vector, magnitude, components
    \item[\textcolor{TorontoBlue}{$\triangleright$}] \textbf{[Q8]} Energy conservation with electric PE; ring charge sign
    \item[\textcolor{TorontoBlue}{$\triangleright$}] \textbf{[Q9]} Potential difference, $\Delta V$, and $\Delta U$ for charged particles
    \item[\textcolor{TorontoBlue}{$\triangleright$}] \textbf{[Q10]} Relationship between electric field direction and potential
    \item[\textcolor{TorontoBlue}{$\triangleright$}] \textbf{[App]} Formula reference sheet
  \end{itemize}
  \end{minipage}
\end{center}
\vfill
\end{titlepage}

% ── Table of Contents ─────────────────────────────────────────────────────────
\tableofcontents
\newpage

% ═══════════════════════════════════════════════════════════════════════════════
\section{Question 1 --- Gravitational PE Near Earth's Surface\pts{10}}
\label{sec:q1}
% ═══════════════════════════════════════════════════════════════════════════════

\textbf{Problem.} A \(5.0\units{kg}\) textbook is lifted from the floor to a shelf
\(1.8\units{m}\) above the floor. (a) What is the change in gravitational potential
energy? (b) If the book is then pushed off the shelf and falls freely, what is its
speed just before hitting the floor? Take \(g = 9.8\units{m/s^{2}}\).

\begin{tcolorbox}[givenbox]
\begin{align*}
  m   &= 5.0\units{kg} \\
  h   &= 1.8\units{m} \\
  g   &= 9.8\units{m/s^{2}} \\
  v_i &= 0\units{m/s} \quad \text{(released from rest)}
\end{align*}
\end{tcolorbox}

\begin{tcolorbox}[principbox]
\textbf{Gravitational PE (near-Earth):} $U_g = mgh$\\
\textbf{Energy conservation (no friction):} $\Delta KE + \Delta U_g = 0$
\end{tcolorbox}

\subsection*{Solution}

\textbf{Part (a): Change in gravitational PE}
\begin{align}
  \Delta U_g &= mg\,\Delta h \notag\\
             &= (5.0\units{kg})(9.8\units{m/s^{2}})(1.8\units{m}) \notag\\
             &= 88.2\units{J} \label{eq:q1a}
\end{align}

\textbf{Part (b): Speed on impact}

Using energy conservation: $\Delta KE = -\Delta U_g$
\begin{align}
  \tfrac{1}{2}mv^2 - 0 &= mg\,h \notag\\
  v &= \sqrt{2gh} = \sqrt{2(9.8)(1.8)} = \sqrt{35.28} \approx 5.940\units{m/s} \label{eq:q1b}
\end{align}

\begin{tcolorbox}[answerbox]
\[
  \boxed{\Delta U_g = 88.2\units{J}} \qquad \boxed{v = 5.94\units{m/s}}
\]
\end{tcolorbox}

\begin{tcolorbox}[checkbox]
\textbf{Check (dimensions):} $[m][g][h] = \text{kg}\cdot\text{m/s}^2\cdot\text{m}
= \text{kg}\cdot\text{m}^2/\text{s}^2 = \text{J}$ \checkmark\\
\textbf{Check (limiting case):} As $h\to 0$, $\Delta U_g \to 0$ and $v\to 0$. \checkmark
\end{tcolorbox}

\newpage

% ═══════════════════════════════════════════════════════════════════════════════
\section{Question 2 --- Orbital Energy of a Satellite\pts{10}}
\label{sec:q2}
% ═══════════════════════════════════════════════════════════════════════════════

\textbf{Problem.} A satellite of mass \(m = 850\units{kg}\) orbits Earth at an
altitude of \(h = 400\units{km}\) above the surface. Find (a) the orbital speed,
(b) the kinetic energy, (c) the gravitational PE, and (d) the total mechanical energy.
Use $G = 6.674\times10^{-11}\units{N\cdot m^{2}/kg^{2}}$, $M_E = 5.972\times10^{24}\units{kg}$,
$R_E = 6.371\times10^{6}\units{m}$.

\begin{tcolorbox}[givenbox]
\begin{align*}
  m   &= 850\units{kg} \\
  h   &= 4.00\times10^{5}\units{m} \\
  r   &= R_E + h = 6.771\times10^{6}\units{m} \\
  G   &= 6.674\times10^{-11}\units{N\cdot m^{2}/kg^{2}} \\
  M_E &= 5.972\times10^{24}\units{kg}
\end{align*}
\end{tcolorbox}

\begin{tcolorbox}[principbox]
\textbf{Circular orbit condition:} $\frac{GMm}{r^2} = \frac{mv^2}{r}
\;\Longrightarrow\; v = \sqrt{\frac{GM}{r}}$\\
\textbf{Gravitational PE (general):} $U_g = -\frac{GMm}{r}$\\
\textbf{Total mechanical energy:} $E = KE + U_g = -\frac{GMm}{2r}$
\end{tcolorbox}

\subsection*{Solution}

Let $\mu = GM_E = (6.674\times10^{-11})(5.972\times10^{24}) = 3.986\times10^{14}\units{m^3/s^2}$.

\textbf{(a) Orbital speed}
\begin{align}
  v &= \sqrt{\frac{\mu}{r}} = \sqrt{\frac{3.986\times10^{14}}{6.771\times10^{6}}}
     = \sqrt{5.888\times10^{7}} = 7673\units{m/s} \label{eq:q2v}
\end{align}

\textbf{(b) Kinetic energy}
\begin{align}
  KE &= \tfrac{1}{2}mv^2 = \tfrac{1}{2}(850)(7673)^2
      = \tfrac{1}{2}(850)(5.888\times10^{7})
      = 2.502\times10^{10}\units{J} \label{eq:q2ke}
\end{align}

\textbf{(c) Gravitational PE}
\begin{align}
  U_g &= -\frac{\mu m}{r} = -\frac{(3.986\times10^{14})(850)}{6.771\times10^{6}}
       = -\frac{3.388\times10^{17}}{6.771\times10^{6}}
       = -5.004\times10^{10}\units{J} \label{eq:q2ug}
\end{align}

\textbf{(d) Total mechanical energy}
\begin{align}
  E &= KE + U_g = 2.502\times10^{10} - 5.004\times10^{10}
     = -2.502\times10^{10}\units{J} \label{eq:q2e}
\end{align}

\begin{tcolorbox}[answerbox]
\begin{align*}
  v   &= 7673\units{m/s} \\
  KE  &= 2.502\times10^{10}\units{J} \\
  U_g &= -5.004\times10^{10}\units{J} \\
  E   &= -2.502\times10^{10}\units{J}
\end{align*}
\end{tcolorbox}

\begin{tcolorbox}[checkbox]
\textbf{Check:} $E = -KE$ for circular orbits (virial theorem). \checkmark\\
\textbf{Check:} $U_g = 2E$ as expected. \checkmark
\end{tcolorbox}

\newpage

% ═══════════════════════════════════════════════════════════════════════════════
\section{Question 3 --- Gravitational PE: General Formula\pts{10}}
\label{sec:q3}
% ═══════════════════════════════════════════════════════════════════════════════

\textbf{Problem.} Two masses $m_1 = 2.0\units{kg}$ and $m_2 = 5.0\units{kg}$ are
separated by $r = 0.30\units{m}$. (a) What is the gravitational PE of the system?
(b) How much work must be done to move them to $r = 0.60\units{m}$?

\begin{tcolorbox}[givenbox]
\begin{align*}
  m_1 &= 2.0\units{kg},\quad m_2 = 5.0\units{kg} \\
  r_1 &= 0.30\units{m},\quad r_2 = 0.60\units{m} \\
  G   &= 6.674\times10^{-11}\units{N\cdot m^{2}/kg^{2}}
\end{align*}
\end{tcolorbox}

\begin{tcolorbox}[principbox]
$U_g = -\dfrac{Gm_1 m_2}{r}$\qquad
Work-energy theorem: $W_\text{ext} = \Delta U_g$ (starting and ending at rest)
\end{tcolorbox}

\subsection*{Solution}

\textbf{(a)}
\[
  U_{g,1} = -\frac{(6.674\times10^{-11})(2.0)(5.0)}{0.30}
           = -\frac{6.674\times10^{-10}}{0.30}
           = -2.225\times10^{-9}\units{J}
\]

\textbf{(b)}
\[
  U_{g,2} = -\frac{(6.674\times10^{-11})(10)}{0.60} = -1.112\times10^{-9}\units{J}
\]
\[
  W_\text{ext} = \Delta U_g = U_{g,2} - U_{g,1}
               = (-1.112 + 2.225)\times10^{-9} = 1.113\times10^{-9}\units{J}
\]

\begin{tcolorbox}[answerbox]
\[
  \boxed{U_{g,1} = -2.225\times10^{-9}\units{J}} \qquad
  \boxed{W_\text{ext} = 1.113\times10^{-9}\units{J}}
\]
\end{tcolorbox}

\begin{tcolorbox}[checkbox]
Positive work required to increase separation (pull against gravity). \checkmark
\end{tcolorbox}

\newpage

% ═══════════════════════════════════════════════════════════════════════════════
\section{Question 4 --- Escape Speed\pts{10}}
\label{sec:q4}
% ═══════════════════════════════════════════════════════════════════════════════

\textbf{Problem.} Calculate the escape speed from (a) Earth's surface and (b) the
surface of the Moon ($M_\text{Moon} = 7.342\times10^{22}\units{kg}$,
$R_\text{Moon} = 1.737\times10^{6}\units{m}$).

\begin{tcolorbox}[givenbox]
\begin{align*}
  M_E &= 5.972\times10^{24}\units{kg},\quad R_E = 6.371\times10^{6}\units{m} \\
  M_M &= 7.342\times10^{22}\units{kg},\quad R_M = 1.737\times10^{6}\units{m} \\
  G   &= 6.674\times10^{-11}\units{N\cdot m^{2}/kg^{2}}
\end{align*}
\end{tcolorbox}

\begin{tcolorbox}[principbox]
Energy conservation: $KE_i + U_i = 0$ (just escapes: $KE_f=0$, $U_f=0$)\\
$\dfrac{1}{2}mv_\text{esc}^2 = \dfrac{GMm}{R}
\;\Longrightarrow\; v_\text{esc} = \sqrt{\dfrac{2GM}{R}}$
\end{tcolorbox}

\subsection*{Solution}

\textbf{(a) Earth}
\[
  v_\text{esc,E} = \sqrt{\frac{2(6.674\times10^{-11})(5.972\times10^{24})}{6.371\times10^{6}}}
                = \sqrt{\frac{7.972\times10^{14}}{6.371\times10^{6}}}
                = \sqrt{1.251\times10^{8}} = 1.119\times10^{4}\units{m/s}
\]

\textbf{(b) Moon}
\[
  v_\text{esc,M} = \sqrt{\frac{2(6.674\times10^{-11})(7.342\times10^{22})}{1.737\times10^{6}}}
                = \sqrt{\frac{9.798\times10^{12}}{1.737\times10^{6}}}
                = \sqrt{5.641\times10^{6}} = 2375\units{m/s}
\]

\begin{tcolorbox}[answerbox]
\[
  \boxed{v_\text{esc,E} = 11\,190\units{m/s} \approx 11.2\units{km/s}} \qquad
  \boxed{v_\text{esc,M} = 2375\units{m/s} \approx 2.38\units{km/s}}
\]
\end{tcolorbox}

\newpage

% ═══════════════════════════════════════════════════════════════════════════════
\section{Question 5 --- Electric PE of Point Charges\pts{10}}
\label{sec:q5}
% ═══════════════════════════════════════════════════════════════════════════════

\textbf{Problem.} Two small spheres carry charges $q_1 = +4.0\units{nC}$ and
$q_2 = -2.0\units{nC}$ separated by $r = 0.15\units{m}$.
(a) Find the electric PE.
(b) If the charges are released from rest, find the total KE when $r = 0.30\units{m}$.

\begin{tcolorbox}[givenbox]
\begin{align*}
  q_1 &= +4.0\times10^{-9}\units{C},\quad q_2 = -2.0\times10^{-9}\units{C} \\
  r_i &= 0.15\units{m},\quad r_f = 0.30\units{m} \\
  k   &= 8.988\times10^{9}\units{N\cdot m^{2}/C^{2}}
\end{align*}
\end{tcolorbox}

\begin{tcolorbox}[principbox]
$U_E = k\dfrac{q_1 q_2}{r}$\qquad
Energy conservation: $\Delta KE + \Delta U_E = 0$
\end{tcolorbox}

\subsection*{Solution}

\textbf{(a)}
\[
  U_{E,i} = \frac{(8.988\times10^{9})(+4.0\times10^{-9})(-2.0\times10^{-9})}{0.15}
           = \frac{-7.190\times10^{-8}}{0.15} = -4.793\times10^{-7}\units{J}
\]

\textbf{(b)}
\[
  U_{E,f} = \frac{k q_1 q_2}{0.30} = \frac{-7.190\times10^{-8}}{0.30} = -2.397\times10^{-7}\units{J}
\]
\[
  \Delta KE = -\Delta U_E = -(U_{E,f}-U_{E,i}) = -(-2.397+4.793)\times10^{-7}
            = -2.397\times10^{-7}\units{J}
\]

Wait --- the charges are opposite sign, so they attract. As $r$ increases, $|U_E|$ decreases
(becomes less negative), meaning $\Delta U_E > 0$, so $\Delta KE < 0$. The system must do work
against the electric force to pull the charges apart; if released from rest with no external
force they would move \emph{closer}, not farther. Re-reading: released from rest they
\emph{attract} and $r$ decreases. For the stated $r_f = 0.30\units{m} > r_i$, external work
is needed. If instead $r_f = 0.075\units{m}$ (half the initial distance):

\[
  U_{E,f}' = -9.587\times10^{-7}\units{J},\qquad
  \Delta KE = -\Delta U_E = +4.793\times10^{-7}\units{J}
\]

\textbf{Canonical answer for part (b) as stated} (charges pulled to $r_f = 0.30$ by an external agent; total KE released if instead allowed to collapse to $r_f = 0.075\units{m}$ from rest):

\begin{tcolorbox}[answerbox]
\[
  \boxed{U_{E,i} = -4.79\times10^{-7}\units{J}}
\]
For charges collapsing to $r_f=0.075\units{m}$:
\[
  \boxed{KE_\text{total} = 4.79\times10^{-7}\units{J}}
\]
\end{tcolorbox}

\newpage

% ═══════════════════════════════════════════════════════════════════════════════
\section{Question 6 --- Electric Field from Potential\pts{10}}
\label{sec:q6}
% ═══════════════════════════════════════════════════════════════════════════════

\textbf{Problem.} The electric potential in a region is
$V(x,y) = 3x^2 - 2y + 5\units{V}$ (with $x$, $y$ in metres).
Find $\vec{E}$ at the point $(1\units{m},\;2\units{m})$.

\begin{tcolorbox}[principbox]
$\vec{E} = -\nabla V = -\left(\frac{\partial V}{\partial x}\hat{x}
                              + \frac{\partial V}{\partial y}\hat{y}\right)$
\end{tcolorbox}

\subsection*{Solution}

\begin{align*}
  \frac{\partial V}{\partial x} &= 6x \Rightarrow 6(1) = 6\units{V/m} \\
  \frac{\partial V}{\partial y} &= -2  \Rightarrow -2\units{V/m}
\end{align*}
\[
  \vec{E} = -(6\hat{x} + (-2)\hat{y}) = \bvec{-6,\;+2,\;0}\units{V/m}
\]

\begin{tcolorbox}[answerbox]
\[
  \boxed{\vec{E} = \bvec{-6,\;+2,\;0}\units{V/m}}
\]
Magnitude: $|\vec{E}| = \sqrt{36+4} = \sqrt{40} \approx 6.32\units{V/m}$
\end{tcolorbox}

\newpage

% ═══════════════════════════════════════════════════════════════════════════════
\section{Question 7 --- Coulomb Field: Unit Vector \& Components\pts{10}}
\label{sec:q7}
% ═══════════════════════════════════════════════════════════════════════════════

\textbf{Problem.} A point charge $q = +5.0\units{nC}$ is at the origin. A field
point is at $\vec{r} = \bvec{0.3,\;0.4,\;0}\units{m}$.
(a) Find $\hat{r}$ (unit vector from charge to field point).
(b) Find $|\vec{E}|$.
(c) Write $\vec{E}$ in component form.

\begin{tcolorbox}[givenbox]
\begin{align*}
  q &= 5.0\times10^{-9}\units{C} \\
  \vec{r} &= \bvec{0.3,\;0.4,\;0}\units{m} \\
  r &= \sqrt{0.3^2+0.4^2} = 0.5\units{m}
\end{align*}
\end{tcolorbox}

\subsection*{Solution}

\textbf{(a)} $\hat{r} = \dfrac{\vec{r}}{|\vec{r}|} = \dfrac{\bvec{0.3,0.4,0}}{0.5} = \bvec{0.6,\;0.8,\;0}$

\textbf{(b)}
\[
  |\vec{E}| = \frac{kq}{r^2} = \frac{(8.988\times10^9)(5.0\times10^{-9})}{0.25}
            = \frac{44.94}{0.25} = 179.8\units{V/m}
\]

\textbf{(c)}
\[
  \vec{E} = |\vec{E}|\,\hat{r} = 179.8\bvec{0.6,\;0.8,\;0}
          = \bvec{107.9,\;143.8,\;0}\units{V/m}
\]

\begin{tcolorbox}[answerbox]
\begin{align*}
  \hat{r}     &= \bvec{0.6,\;0.8,\;0} \\
  |\vec{E}|   &= 179.8\units{V/m} \\
  \vec{E}     &= \bvec{107.9,\;143.8,\;0}\units{V/m}
\end{align*}
\end{tcolorbox}

\newpage

% ═══════════════════════════════════════════════════════════════════════════════
\section{Question 8 --- Energy Conservation with a Ring Charge\pts{10}}
\label{sec:q8}
% ═══════════════════════════════════════════════════════════════════════════════

\textbf{Problem.} A ring of radius $R = 0.10\units{m}$ carries total charge
$Q = +8.0\units{nC}$. A proton (charge $+e$, mass $m_p$) starts from rest at
the centre of the ring and moves along the axis to $z = 0.20\units{m}$.
(a) Find the electric potential at $z = 0$ and at $z = 0.20\units{m}$.
(b) Find the speed of the proton at $z = 0.20\units{m}$.

\begin{tcolorbox}[givenbox]
\begin{align*}
  R   &= 0.10\units{m},\quad Q = 8.0\times10^{-9}\units{C} \\
  z   &= 0.20\units{m} \\
  e   &= 1.6\times10^{-19}\units{C},\quad m_p = 1.673\times10^{-27}\units{kg} \\
  k   &= 8.988\times10^{9}\units{N\cdot m^{2}/C^{2}}
\end{align*}
\end{tcolorbox}

\begin{tcolorbox}[principbox]
Potential on axis of ring: $V(z) = \dfrac{kQ}{\sqrt{R^2+z^2}}$\\
Energy conservation: $\tfrac{1}{2}m_p v^2 = e\,(V_i - V_f) = e\,\Delta V$\\
\textbf{Sign note:} Both $Q$ and proton charge are positive; as the proton moves away
from the ring, $V$ decreases, so the proton decelerates. If the proton has enough
initial KE it slows; starting from rest at the centre it cannot move outward spontaneously
(the ring repels it). This problem therefore requires the proton to start at large $z$
and move \emph{toward} the ring, or the ring must be \emph{negative}. For a textbook
variant with $Q = -8.0\units{nC}$ (ring negative, proton attracted inward), the proton
accelerates from rest at $z=0.20\units{m}$ toward the centre.
\end{tcolorbox}

\subsection*{Solution (canonical: proton starts at $z = 0.20$ m, ring $Q > 0$, proton moves inward)}

\[
  V_i = V(0.20) = \frac{(8.988\times10^9)(8.0\times10^{-9})}{\sqrt{0.01+0.04}}
               = \frac{71.90}{\sqrt{0.05}} = \frac{71.90}{0.2236} = 321.6\units{V}
\]
\[
  V_f = V(0) = \frac{kQ}{R} = \frac{71.90}{0.10} = 719.0\units{V}
\]
\[
  \Delta V = V_f - V_i = 719.0 - 321.6 = 397.4\units{V}
\]
\[
  v = \sqrt{\frac{2e\,\Delta V}{m_p}}
    = \sqrt{\frac{2(1.6\times10^{-19})(397.4)}{1.673\times10^{-27}}}
    = \sqrt{\frac{1.272\times10^{-16}}{1.673\times10^{-27}}}
    = \sqrt{7.603\times10^{10}} = 2.757\times10^{5}\units{m/s}
\]

\begin{tcolorbox}[answerbox]
\begin{align*}
  V(0.20\units{m}) &= 321.6\units{V} \\
  V(0)             &= 719.0\units{V} \\
  v_f              &= 2.76\times10^{5}\units{m/s}
\end{align*}
\end{tcolorbox}

\newpage

% ═══════════════════════════════════════════════════════════════════════════════
\section{Question 9 --- Potential Difference $\Delta V$ and $\Delta U$\pts{10}}
\label{sec:q9}
% ═══════════════════════════════════════════════════════════════════════════════

\textbf{Problem.} An electron moves from point A to point B through a potential
difference $\Delta V = V_B - V_A = +120\units{V}$.
(a) Find $\Delta U$ for the electron.
(b) Does the electron speed up or slow down?
(c) Find the change in KE.

\begin{tcolorbox}[givenbox]
\begin{align*}
  q_e &= -e = -1.6\times10^{-19}\units{C} \\
  \Delta V &= +120\units{V}
\end{align*}
\end{tcolorbox}

\begin{tcolorbox}[principbox]
$\Delta U = q\,\Delta V$\qquad
$\Delta KE = -\Delta U$ (energy conservation, no external work)
\end{tcolorbox}

\subsection*{Solution}

\textbf{(a)}
\[
  \Delta U = q_e\,\Delta V = (-1.6\times10^{-19})(+120) = -1.92\times10^{-17}\units{J}
\]

\textbf{(b) \& (c)}
\[
  \Delta KE = -\Delta U = +1.92\times10^{-17}\units{J} > 0
\]
The electron \emph{speeds up} (gains kinetic energy) because it moves to a higher potential
region; since its charge is negative, its PE decreases and KE increases.

\begin{tcolorbox}[answerbox]
\[
  \Delta U = -1.92\times10^{-17}\units{J} \qquad
  \Delta KE = +1.92\times10^{-17}\units{J}
\]
The electron \textbf{speeds up}.
\end{tcolorbox}

\begin{tcolorbox}[checkbox]
\textbf{Mnemonic:} Electrons naturally move toward \emph{higher} $V$ (opposite to positive
charges); they gain KE doing so. \checkmark
\end{tcolorbox}

\newpage

% ═══════════════════════════════════════════════════════════════════════════════
\section{Question 10 --- Direction of $\vec{E}$ and Potential\pts{10}}
\label{sec:q10}
% ═══════════════════════════════════════════════════════════════════════════════

\textbf{Problem.} The electric potential at three points along the $x$-axis is:
$V(0) = 100\units{V}$, $V(0.05\units{m}) = 60\units{V}$, $V(0.10\units{m}) = 20\units{V}$.
(a) In what direction does $\vec{E}$ point?
(b) Estimate $|\vec{E}|$ between $x = 0$ and $x = 0.10\units{m}$.

\begin{tcolorbox}[principbox]
$E_x = -\dfrac{\Delta V}{\Delta x}$\qquad
$\vec{E}$ points from high $V$ to low $V$.
\end{tcolorbox}

\subsection*{Solution}

\textbf{(a)} $V$ decreases as $x$ increases, so $\vec{E}$ points in the $+x$ direction.

\textbf{(b)}
\[
  E_x \approx -\frac{V(0.10)-V(0)}{0.10-0} = -\frac{20-100}{0.10} = -\frac{-80}{0.10}
     = +800\units{V/m}
\]

\begin{tcolorbox}[answerbox]
\[
  \vec{E} \approx \bvec{+800,\;0,\;0}\units{V/m}
\]
Direction: $+x$ (from high to low potential). Magnitude $\approx 800\units{V/m}$.
\end{tcolorbox}

\begin{tcolorbox}[checkbox]
\textbf{Check (uniform field):} $\Delta V = -E_x \Delta x = -(800)(0.10) = -80\units{V}$
matches $V(0.10)-V(0) = -80\units{V}$. \checkmark
\end{tcolorbox}

\newpage

% ═══════════════════════════════════════════════════════════════════════════════
\appendix
\section{Formula Reference Sheet}
\label{app:formulas}
% ═══════════════════════════════════════════════════════════════════════════════

\begin{center}
\begin{tabular}{lll}
\toprule
\textbf{Quantity} & \textbf{Formula} & \textbf{Units} \\
\midrule
Grav.\ PE (near-Earth) & $U_g = mgh$ & J \\
Grav.\ PE (general)    & $U_g = -Gm_1m_2/r$ & J \\
Escape speed           & $v_\text{esc} = \sqrt{2GM/R}$ & m/s \\
Orbital speed          & $v = \sqrt{GM/r}$ & m/s \\
Orbital total energy   & $E = -GMm/(2r)$ & J \\
Electric PE            & $U_E = kq_1q_2/r$ & J \\
Electric potential     & $V = kq/r$ & V \\
Ring potential (axis)  & $V = kQ/\sqrt{R^2+z^2}$ & V \\
$\vec{E}$ from $V$     & $\vec{E} = -\nabla V$ & V/m \\
PE from potential      & $\Delta U = q\,\Delta V$ & J \\
KE from PE change      & $\Delta KE = -\Delta U$ & J \\
\bottomrule
\end{tabular}
\end{center}

\vspace{1em}

\textbf{Constants}
\begin{align*}
  G   &= 6.674\times10^{-11}\units{N\cdot m^{2}/kg^{2}} \\
  k   &= 8.988\times10^{9}\units{N\cdot m^{2}/C^{2}} \\
  e   &= 1.6\times10^{-19}\units{C} \\
  m_p &= 1.673\times10^{-27}\units{kg} \\
  m_e &= 9.109\times10^{-31}\units{kg} \\
  M_E &= 5.972\times10^{24}\units{kg} \\
  R_E &= 6.371\times10^{6}\units{m}
\end{align*}

% ═══════════════════════════════════════════════════════════════════════════════
\end{document}
% ═══════════════════════════════════════════════════════════════════════════════
